\documentclass[14pt]{mmcs_article}

\usepackage[T2A]{fontenc}
\usepackage[utf8]{inputenc}
\usepackage[russian]{babel}
\usepackage{amsmath,amssymb}
\usepackage{booktabs}
\usepackage{multirow}
\usepackage{url}
\usepackage{svg}
\graphicspath{{images/}}
\svgpath{images/}
\setsvg{inkscapelatex=false}

\newcommand{\rem}{\mathrm{REM}}
\newcommand{\elpd}{\mathrm{elpd}}

\begin{document}

\begin{center}
{\Large\bfseries Определение стресса крысы в окрестности\\
сейсмического события байесовским методом}
\end{center}

\begin{abstract}
Исследуется гипотеза о физиологическом отклике крыс на процессы,
сопровождающие подготовку сейсмического события. В качестве косвенного
индикатора стрессовой реакции рассматривается динамика доли
парадоксального сна (REM) по данным непрерывной регистрации
биоэлектрической активности мозга. Выборка включает 14 дат сейсмических
событий 2022--2025 гг. и 20 пар ``событие--крыса'' с доступной
гипнограммой за восемь суток до толчка. Для выявления смены режима
REM-профиля построена байесовская модель одной точки разладки с
маргинализацией её дискретного положения и оцениванием методом NUTS.
Выполнен полный перебор 896 конфигураций предобработки и
правдоподобий; модели сопоставлялись по PSIS-LOO ожидаемой
логарифмической предсказательной плотности. Лучшие конфигурации
используют среднюю суточную долю REM, распределение Стьюдента и
двухчасовое окно. Для лучшего квартиля моделей апостериорное среднее
положения разладки составляет $6{,}64\pm0{,}42$ суток до события.
Отсутствие влиятельных наблюдений в топ-10 моделей подтверждает
устойчивость результата. Результат следует интерпретировать как
свидетельство изменения REM-ассоциированного физиологического
состояния, а не как прямое измерение стресса или средство прогноза
землетрясений.

\medskip
\noindent\textbf{Ключевые слова:} парадоксальный сон, крыса,
сейсмическое событие, точка разладки, байесовский вывод, PSIS-LOO.
\end{abstract}

\section{Введение}

Точный физический механизм подготовки землетрясений пока не установлен,
однако для сейсмически активных регионов описаны геофизические и
геохимические предвестники \cite{Rodkin2011,Cicerone2009,Conti2021}.
В литературе также накоплены наблюдения необычного поведения животных
перед толчками; возможная биологическая основа такой чувствительности
обсуждается, в частности, в работах
\cite{Kirschvink2000,Grant2011,Yokoi2003}. Эти сведения сами по себе
не доказывают наличие универсального предвестника и требуют проверки на
инструментально зарегистрированных физиологических данных.

Стресс и тревога изменяют цикл сон--бодрствование и структуру
парадоксального сна у грызунов
\cite{Rampin1991,Meerlo1997,Palma2000,Adrien1991,Sanford2010}.
Парадоксальный сон характеризуется специфической электроэнцефалографической
картиной и у крыс входит в состав выраженного суточного ритма
\cite{Vyazovskiy2005,Franken1991,WeberDan2016}. Поэтому непрерывная
регистрация сна позволяет проверить, меняется ли физиологическая динамика
вблизи сейсмического события.

Особенность рассматриваемой задачи состоит в сверхмалой выборке и
неоднородности наблюдений. Редкие события, индивидуальные различия
животных и автокорреляция временных рядов ограничивают применимость
сложных ``чёрных ящиков'' и требуют явного учёта неопределённости
\cite{Hyndman2018,Sinkovec2022}. Цель работы --- разработать
интерпретируемый байесовский метод поиска смены режима в динамике REM и
оценить время такой смены относительно основного толчка. Работа не
направлена на оперативный прогноз землетрясений: проверяется только
гипотеза о статистически различимом физиологическом отклике.

\section{Материалы и методы}

\subsection{Экспериментальные данные}

Электрокортикограмма и электрическая активность гиппокампа непрерывно
регистрировались у взрослых самцов крыс линии пасюк на станции
``Институт'' Камчатского филиала ЕГС РАН (Петропавловск-Камчатский).
Животных содержали индивидуально при температуре $23\pm1^{\circ}$C и
режиме освещения 12 ч свет / 12 ч темнота. Электроды в области
сенсомоторной коры и поле CA1 гиппокампа соединялись с 32-канальной
системой регистрации; частота оцифровки составляла 1 кГц. Все процедуры
были одобрены комиссией по биоэтике Южного федерального университета.

Для анализа отобраны 14 дат землетрясений 2022--2025 гг. с
инструментальной интенсивностью более двух баллов и доступными записями
за восемь суток до события. Две даты 2025 года представлены записями
четырёх животных каждая; поэтому фактическая выборка состоит из 20
пар ``событие--крыса''. Такая формулировка отделяет число уникальных
сейсмических дат от числа физиологических реализаций.

\subsection{Получение REM-профилей}

Суточную запись разбивали на неперекрывающиеся эпохи длительностью 5 с.
После исключения артефактов для каждой эпохи вычисляли спектры мощности.
Медленный сон определяли по высокой мощности ЭКоГ в
дельта-диапазоне 2,5--5,5 Гц, а парадоксальный сон --- по высоким
значениям отношения тета- к дельта-мощности в сигнале гиппокампа
\cite{Fang2010,Vyazovskiy2005,Saevskiy2025}. Результатом была
гипнограмма с метками бодрствования, медленного и парадоксального сна
(рис.~\ref{fig:hypnogram}).

\begin{figure}[H]
  \centering
  \includesvg[width=\textwidth]{hypno.svg}
  \caption{Пример гипнограммы крысы: последовательность состояний
  бодрствования, медленного и парадоксального сна, полученная по
  непрерывной записи биоэлектрической активности.}
  \label{fig:hypnogram}
\end{figure}

Из гипнограммы строили суточный REM-профиль --- ряд долей времени,
проведённого в REM, в скользящих окнах шириной $W$ и шагом $S$. Для
каждого профиля рассчитывали среднее
\[
 \bar P=\frac{1}{N}\sum_{i=1}^{N}p_i
\]
и размах $R=\max(P)-\min(P)$. Признаки вычисляли для объединённых
дневного и ночного фрагментов (\texttt{concat}), а также раздельно для
чётных и нечётных временных чанков. Каждый восьмидневный профиль
масштабировали в диапазон $[0,1]$ независимо для каждой пары
``событие--крыса'', чтобы анализировать форму динамики, а не абсолютный
уровень REM.

\subsection{Байесовская модель точки разладки}

Для каждого набора признаков использовали модель с одной точкой
разладки $\tau$, разделяющей восемь дней до толчка на два режима:
\[
 \mathbf{x}_t\sim
 \begin{cases}
 p(\mathbf{x}_t\mid\boldsymbol{\theta}_1),&t<\tau,\\
 p(\mathbf{x}_t\mid\boldsymbol{\theta}_2),&t\geq\tau.
 \end{cases}
\]
Положение $\tau$ задавали дискретной равномерной априорной величиной,
а параметры $\boldsymbol{\theta}_1$ и $\boldsymbol{\theta}_2$ снабжали
слабоинформативными априорными распределениями. Дискретную переменную
маргинализировали из целевой плотности; это сделало возможным применение
градиентного сэмплера NUTS \cite{HoffmanGelman2014,
lai2023automaticallymarginalizedmcmcprobabilistic}. После
сэмплирования распределение $p(\tau\mid\mathcal D)$ восстанавливали
аналитическим усреднением условных вероятностей положений разладки.

\begin{figure}[H]
  \centering
  \includesvg[width=\textwidth]{grafviz.svg}
  \caption{Вероятностная структура модели с общей точкой разладки:
  параметры распределений REM-признаков различаются до и после $\tau$.}
  \label{fig:model}
\end{figure}

\subsection{Полный перебор и сравнение моделей}

В пространстве поиска варьировали ширину окна
$W\in\{2,3,6\}$ ч, шаг $S\in\{1,2\}$ ч, наборы признаков
\texttt{mean} и \texttt{range} в группировках \texttt{concat},
\texttt{even}, \texttt{odd}, а также семейства правдоподобий
\texttt{student\_t}, \texttt{lognormal} и \texttt{beta}.
Всего было оценено 896 конфигураций. Каждая модель сэмплировалась в
четырёх независимых цепях; использовали 3000 итераций разогрева и
1000 сохраняемых итераций на цепь.

Качество обобщения оценивали методом PSIS-LOO
\cite{Vehtari2017,Kumar2019}. В качестве наблюдения в LOO выступала
целая пара ``событие--крыса'', поэтому сравнивали нормированную метрику
\[
\overline{\overline{\elpd}}_{\mathrm{LOO}}=
\frac{\widehat{\elpd}_{\mathrm{LOO}}}{F E'},
\]
где $F$ --- число активных признаков, $E'$ --- число записей,
оставшихся после возможного исключения влиятельных наблюдений. Контроль
качества включал $\hat R\leq1{,}05$, достаточный эффективный размер
выборки, отсутствие дивергенций и диагностику Парето $\hat k$.

\section{Результаты}

\subsection{Качество и структура пространства моделей}

Все 896 конфигураций прошли заданные критерии сходимости: $\hat R$ не
превышал 1,05, эффективный размер выборки был не менее 500, дивергенции
NUTS отсутствовали. Распределение нормированного ELPD имело широкое
плато для основной массы конфигураций и выделенный верхний хвост
моделей-лидеров. В центральной части оно близко к нормальному
распределению с параметрами $\mu=3{,}20$, $\sigma=1{,}01$; для
Q--Q-аппроксимации получено $R^2=0{,}975$.

Результаты полного перебора показывают, что информативным является
прежде всего суточное среднее REM. Признак \texttt{concat:mean}
встречается в 94\% моделей лучшего квартиля и отсутствует в худшем
квартиле; все 45 моделей из лучших 5\% содержат этот признак
(точный критерий Фишера, $p<0{,}001$, отношение шансов $2{,}0$).
Напротив, раздельный анализ дневных и ночных частей, а также добавление
размаха не увеличивали предсказательную способность.

\begin{table}[H]
\centering
\caption{Распространённость REM-признаков в квартилях качества моделей,
\%. Q1 --- лучший квартиль по нормированному ELPD.}
\label{tab:features}
\begin{tabular}{lrrrr}
\toprule
Признак & Q1 & Q2 & Q3 & Q4\\
\midrule
\texttt{concat:mean}  & 94 & 80 & 40 & 0\\
\texttt{concat:range} & 53 & 54 & 62 & 45\\
\texttt{even:mean}    & 47 & 54 & 63 & 50\\
\texttt{odd:mean}     & 43 & 55 & 62 & 54\\
\texttt{even:range}   & 27 & 58 & 60 & 69\\
\texttt{odd:range}    & 38 & 55 & 60 & 62\\
\bottomrule
\end{tabular}
\end{table}

\subsection{Оптимальная конфигурация}

Наиболее предпочтительный класс моделей использует
\texttt{concat:mean} с распределением Стьюдента. Это согласуется с
распределением семейств правдоподобия: в первом квартиле
\texttt{student\_t} выбран для среднего в 64\% моделей, тогда как в
четвёртом квартиле --- только в 27\%. Двухчасовое окно обеспечивает
средний ELPD $3{,}46\pm1{,}07$, существенно превышающий
$2{,}65\pm0{,}72$ для окна 6 ч. Шаг окна практически не влияет на
результат: его средние значения ELPD составляют 3,18 и 3,23 для
$S=1$ и $S=2$ ч соответственно (табл.~\ref{tab:window}).

\begin{table}[H]
\centering
\caption{Средний нормированный ELPD в зависимости от параметров
формирования REM-профиля.}
\label{tab:window}
\begin{tabular}{llr}
\toprule
Параметр & Значение & $\overline{\elpd}_{\mathrm{LOO}}\pm\sigma$\\
\midrule
\multirow{3}{*}{Ширина окна $W$} & 2 ч & $3{,}46\pm1{,}07$\\
 & 3 ч & $3{,}21\pm0{,}91$\\
 & 6 ч & $2{,}65\pm0{,}72$\\
\midrule
\multirow{2}{*}{Шаг $S$} & 1 ч & $3{,}18\pm1{,}01$\\
 & 2 ч & $3{,}23\pm1{,}02$\\
\bottomrule
\end{tabular}
\end{table}

Три лидирующие конфигурации статистически неразличимы при попарном
сравнении ($p>0{,}05$); они образуют класс практически эквивалентных
решений. Начиная с четвёртого ранга различия с лучшей моделью значимы
($p<0{,}05$), а для моделей с седьмого по десятое место
$p<0{,}001$. Добавление \texttt{range} сужает апостериорное
распределение $\tau$ в $1{,}3$--$1{,}9$ раза, но ухудшает ELPD:
разница с лучшими моделями без этого признака составляет 52,5 единицы
($p<0{,}001$). Тем самым точность локализации и предсказательная
способность образуют компромисс, в котором для основной модели
предпочтительно суточное среднее REM.

\subsection{Локализация изменения REM-динамики}

Апостериорное положение точки разладки устойчиво к выбору качества
модели (табл.~\ref{tab:tau}). В лучшем квартиле среднее
$\mathbb E[\tau]$ равно $6{,}64\pm0{,}42$ суток, а медиана --- 6,27
суток. Даже между лучшим и худшим квартилями разность средних оценок
составляет только 0,18 суток. Таким образом, согласованная оценка
помещает смену режима REM-профиля на 6--7-е сутки до основного толчка.
Именно эта оценка используется в статье; упоминание окна 2--3 суток в
одном из исходных фрагментов диплома не согласуется с его таблицей
результатов и итоговыми выводами.

\begin{table}[H]
\centering
\caption{Апостериорные характеристики $\tau$ по квартилям ELPD.}
\label{tab:tau}
\begin{tabular}{lrrrr}
\toprule
Квартиль & Диапазон ELPD & $\mathbb E[\tau]$ & $Q_2(\tau)$ &
$\mathrm{std}(\tau)$\\
\midrule
Q1 (лучшие 25\%) & $[3{,}84;7{,}97]$ & $6{,}64\pm0{,}42$ & 6,27 & 0,88\\
Q2 & $[3{,}15;3{,}83]$ & $6{,}71\pm0{,}56$ & 6,32 & 0,71\\
Q3 & $[2{,}45;3{,}15]$ & $6{,}54\pm0{,}70$ & 6,15 & 0,81\\
Q4 (худшие 25\%) & $[0{,}15;2{,}44]$ & $6{,}46\pm0{,}88$ & 6,19 & 1,00\\
\bottomrule
\end{tabular}
\end{table}

Для топ-10 моделей не выявлены записи с $\hat{k}\geq0{,}7$; повторное
обучение после исключения влиятельных событий не потребовалось, а
$\Delta\elpd=0$. Следовательно, высокие оценки моделей-лидеров не
обусловлены одной экстремальной записью. Это поддерживает устойчивость
обнаруженной точки разладки в рамках доступной выборки, но не отменяет
необходимости независимой проверки на новых данных.

\section{Обсуждение}

Полученный результат согласуется с предположением, что перед частью
сейсмических событий меняется физиологическая динамика крыс, отражаемая
в среднем уровне REM. Байесовская постановка особенно полезна здесь:
вместо единственной даты переключения она даёт полное распределение
неопределённости положения $\tau$, а PSIS-LOO позволяет сравнить
большое число моделей без повторного обучения для каждого исключаемого
наблюдения \cite{Vehtari2017}.

Вместе с тем результат нельзя интерпретировать как прямое обнаружение
стресса. В исследовании не измерялись кортикостерон, поведенческие
показатели тревоги и другие независимые биомаркеры. REM выступает
косвенным физиологическим индикатором, который может реагировать на
несколько факторы. Кроме того, выборка мала, содержит записи разных
животных и событий, а модель допускает лишь единственную скачкообразную
смену режима. Более сложные варианты --- плавный переход, несколько
точек разладки и иерархическая структура для животных и характеристик
землетрясений --- требуют расширения данных.

Перспективным направлением является проспективная проверка заранее
зафиксированной конфигурации \texttt{concat:mean} с двухчасовым окном
на новых сейсмических событиях. Необходимо также сопоставить силу
эффекта с магнитудой, глубиной и расстоянием до эпицентра, а
физиологическую интерпретацию --- подтвердить независимыми маркерами
стресса.

\section{Выводы}

\begin{enumerate}
  \item Разработан метод анализа восьмидневных REM-профилей на основе
  байесовской модели одной точки разладки с маргинализацией дискретного
  параметра $\tau$.
  \item В полном переборе 896 конфигураций наилучший класс моделей
  использует суточное среднее REM (\texttt{concat:mean}),
  распределение Стьюдента и окно шириной 2 ч.
  \item Оценка положения разладки устойчива: для лучших моделей
  $\mathbb E[\tau]=6{,}64\pm0{,}42$ суток до сейсмического события,
  а между квартилями качества разность средних не превышает 0,2 суток.
  \item Топ-10 моделей не содержат влиятельных записей по критерию
  Парето $\hat{k}\geq0{,}7$, что поддерживает робастность результата
  внутри исследуемой выборки.
  \item Полученные данные указывают на REM-ассоциированное изменение
  физиологического состояния вблизи сейсмических событий; для вывода о
  стрессе и переносимости результата нужны независимые биомаркеры и
  проспективная валидация.
\end{enumerate}

\renewcommand{\refname}{\centering\textbf{Литература}}
\bibliographystyle{ugost2008-local}
\bibliography{mmcs-thesis2020}

\end{document}
